% Appendix G, the natural arm's screening, replay and construction (restructure plan, part 3). Every number is a macro:
% \nat (generated/natural/numbers.tex; evidence_ledger_natural.md: data/manifests/screening_natural_aidev*.json,
% natural_tasks_index.json, rejection_ledger.json), \dv (evidence_ledger_design.md) or \fv; the per-family tables are
% make_tables.py --natural-run's (evidence_ledger_natural_run.md). Moved here: data.tex's populations and natural arm,
% results_natural.tex's funnel and replay, and the appendix's natural tables; the comparison with the controlled study is
% the table the natural section points to. Families are named, never counted, where a statement concerns the run's
% families; S9 and S10 counts are screening counts. No thread text, repository name, login or agent product is written.
\section{Natural arm: screening, replay and construction}
\label{app:natural}\label{sec:natural}\label{sec:results-natural}

\paragraph{Populations and sources.} We keep populations of these kinds apart and never pool them without saying so
(protocol~\S2, \S15.2): \emph{natural} data from real use of coding agents, a population per source; \emph{controlled}
sessions, real agent runs whose user turns the study's agent wrote beforehand, the same model as the session agent, with
no person writing or reviewing them; and synthetic \emph{fixtures} for software tests. SWE-chat \cite{swechat} was not
used: its access terms were not accepted (D4). Ungated, openly licensed sources were screened instead (D47): Trace
Commons \cite{tracecommons2026trace}, whose agent sessions were all rejected (\nat{tc.rejected} of \nat{tc.screened}),
none of them reconstructable (Table~\ref{tab:trace-commons}); and AIDev v5 \cite{li2025rise,li2026aidev,haoli2025aidev}, which
supplies the natural arm. Every verifier result outside the natural case studies is on the controlled population, whose
feedback is synthetic.

\paragraph{The source.} In the natural arm a reviewer's comment on an agent's pull request (PR) takes the place of the
user's correction. AIDev v5 is a CC-BY-4.0 compilation of agent-authored PRs in which each repository keeps its own
licence; its files are pinned to the tag recorded in D47 and hash-checked.

\paragraph{The deterministic funnel (S1 to S8a).} The stages keep (protocol~\S15.6; Table~\ref{tab:natural-funnel}):
\begin{itemize}
\item PRs of Python repositories under MIT, Apache-2.0 or BSD licences that are not forks;
\item thread-starting inline comments by a person other than the PR's author, on a \texttt{.py} file with a diff hunk;
\item whose commented code changed later, with the reviewed commit in the PR's commit list and a later commit;
\item on PRs linked to an issue written by a person;
\item whose first later commit is attributed to an agent.
\end{itemize}
From \nat{aidev.s1.prs} agent-authored PRs in \nat{aidev.s1.repos} repositories, \nat{aidev.s8a.comments} comments
survive, on \nat{aidev.s8a.prs} PRs in \nat{aidev.s8a.repos} repositories. Attribution needs a bot account or an agent
co-author trailer, which misses agents committing as a person, and the most frequent agent product accounts for
\nat{aidev.s8a.topagent} of the \nat{aidev.s8a.comments} comments (D48).

\paragraph{The model-judged stage (S9).}
\begin{itemize}
\item Each sampled comment goes to \dv{screen.s9.judges.w} independent configuration-A judges, who read the pseudonymised
  comment, its diff hunk, the PR's title and the linked issue, and the diff from the reviewed commit to the first later commit that touches the commented
  file (partial clones, at most \dv{screen.s9.repo}\,MB per repository and \dv{screen.s9.total}\,GB in all).
\item A comment qualifies only if both find that it asks for a behaviour change or reports a defect, that this commit
  addresses it, and that an offline Python test can check it; disagreements are excluded. These verdicts are model
  proposals, queued for human coding.
\item The \emph{first sample} is \nat{s9.first.comments} of the \nat{aidev.s8a.comments} comments, drawn with a fixed
  seed (D47, D48); \dv{screen.s9.rejudged} comments first judged on an empty diff were judged again (D50). In it,
  \nat{s9.first.qualifying} comments qualify (\nat{s9.first.rate}, Wilson \nat{s9.first.wilson}).
\item The \emph{extension} is the other \nat{s9.ext.comments}, screened in comment-id order under unchanged rules by new
  judges blind to the first sample's judgments (protocol~\S19.2); \nat{s9.ext.qualifying} qualify (\nat{s9.ext.rate},
  Wilson \nat{s9.ext.wilson}). It was decided after the first sample yielded \nat{s10.first.passed} replayable tasks,
  short of the \dv{design.size.threshold} below which protocol~\S15.6 reports the arm as a screening study only, and
  before any natural contract, verifier, held-out candidate or label existed: a post hoc design change (D60).
\item The leading reason a comment fails is that it asks for no behaviour change or defect fix
  (\nat{s9.first.fail.behaviour} comments of the first sample and \nat{s9.ext.fail.behaviour} of the extension). The shares are model-judged screening
  yields, not prevalence.
\end{itemize}

\IfFileExists{generated/natural/natural_funnel.tex}{\begin{table}[t]\centering\footnotesize
\caption{Screening funnel of the natural PR-review arm (AIDev v5). Top: the deterministic stages S1--S8a; S1 and S2 count PRs. Without the linked-issue requirement of S8, 3{,}126 comments would pass S8 and 2{,}263 would pass S8a. Bottom: S9 for the first sample (a random sample of the S8a comments) and for the extension (every other S8a comment, decided after the first sample's S10 yield was known, protocol \S19), then families, S10 (replayability) and S11. \textbf{The S9 verdicts are model proposals of two judges, not human coding.} In the first sample the comments without reachable commits were judged without a diff and failed as not fetchable; in the extension they failed before judging. Extension, share of the judged comments: 22.5\% (18.7\%--26.8\%). First sample, judge $\kappa$ over the 131 comments with reachable commits: 0.93; 0.91; 0.93. No natural verifier, label or run exists yet.}\label{tab:natural-funnel}
\begin{tabular}{llrrr}\toprule Stage & Rule & Comments & PRs & Repositories \\\midrule
S1 & Agent-authored pull requests (AIDev v5) & -- & 361{,}296 & 25{,}402 \\
S2 & Python repository; MIT, Apache-2.0 or BSD licence; not a fork & -- & 44{,}914 & 3{,}333 \\
S3 & Inline review comments on these PRs & 28{,}732 & 5{,}030 & 757 \\
S4 & Written by a person, starts a thread, not by the PR's author & 6{,}751 & 2{,}070 & 427 \\
S5 & On a \texttt{.py} file, with a diff hunk & 3{,}895 & 1{,}241 & 306 \\
S6 & Commented code changed afterwards (a proxy) & 3{,}540 & 1{,}085 & 268 \\
S7 & Reviewed commit in the PR's commit list; a later commit exists & 3{,}126 & 987 & 246 \\
S8 & PR linked to an issue written by a person & 926 & 255 & 92 \\
S8a & First later commit attributed to an agent & 684 & 190 & 63 \\
\bottomrule\end{tabular}

\medskip
\begin{tabular}{p{0.42\linewidth}p{0.22\linewidth}p{0.22\linewidth}}\toprule S9 to S11 & First sample & Extension \\\midrule
Comments screened & 150 & 534 \\
\quad commits reachable (diff computable) & 131 & 419 \\
\quad clone refused by the session's proxy & 17 & 56 \\
\quad repository over the fetch cap & 2 & 59 \\
Judged by both judges & 150 & 405 \\
\quad on a non-empty diff & 130 & 405 \\
Failed before judging & 0 & 129 \\
Qualifying (share of screened; Wilson 95\%) & 31 (20.7\%; 15.0\%--27.8\%) & 91 (17.0\%; 14.1\%--20.5\%) \\
\quad share of comments with reachable commits & 23.7\% (17.2\%--31.6\%) & 21.7\% (18.0\%--25.9\%) \\
Excluded: the judges disagree & 2 & 10 \\
Failed (one primary reason each) & 117 & 433 \\
\quad repository over the fetch cap & -- & 59 \\
\quad commits not fetchable & 19 & 56 \\
\quad no later commit changes the commented file & -- & 14 \\
\quad addressing commit not agent-attributed & 0 & 1 \\
\quad not a behaviour change or defect & 92 & 278 \\
\quad not addressed by the later commit & 6 & 23 \\
\quad not testable offline & 0 & 2 \\
Qualifying PRs / repositories & 25 / 14 & 42 / 19 \\
Judge $\kappa$: behaviour or defect; addressed; testable offline & 0.93; 0.94; 0.93 (n = 150) & 0.95; 0.89; 0.94 (n = 405) \\\midrule
Families + reserves (at most 2 per repository) & 20 + 5 & 19 + 4 \\
S10 processed / passed / rejected & 22 / 6 / 16 & \parbox[t]{\linewidth}{\raggedright \unfilled{pending: S10 of the S9 extension}} \\
\quad rejected: base tests fail to collect & 8 & \unfilled{pending} \\
\quad rejected: no resolvable dependency closure & 4 & \unfilled{pending} \\
\quad rejected: environment too large & 1 & \unfilled{pending} \\
\quad rejected: empty base-test subset & 1 & \unfilled{pending} \\
\quad rejected: fetch over the repository cap & 1 & \unfilled{pending} \\
\quad rejected: merge rule & 1 & \unfilled{pending} \\
S11 accepted for building & 6 & \unfilled{pending} \\
\bottomrule\end{tabular}\end{table}
}{\unfilled{natural funnel table}}

\paragraph{Families and replay (S10).}
\begin{itemize}
\item A family is a PR, corrected by its earliest qualifying comment. At most \dv{natural.family.repo.cap.w} come from a
  repository, the earliest by creation time that pass S10, and a rejected family is replaced by its repository's next
  qualifying PR (protocol~\S17.3). The first sample gives NP-01 to NP-20 and reserves NP-21 to NP-25; the extension's
  families, from NP-26, only fill open places, and first-sample families keep the corrections they were built with, a
  builder default that departs from the earliest-comment rule where an extension comment came first, as for NP-05 (D63,
  D66; Appendix~\ref{app:deviations}).
\item S10, fixed before any build (protocol~\S15.6, \S17), requires reachable reviewed and addressing commits; a base by
  a fixed merge-base rule, approximate where the rule falls back or a merge leaves the commented file unchanged, and
  rejected where a merge changes it; a pinned, wheels-only Python~3.12 environment of at most
  \dv{screen.s10.wheelhouse}\,GB, resolved from the dependencies declared at the base plus the packages the PR adds
  (pure-Python sdists may be built offline); and a base-test subset (test files that import a changed module or that the
  PR changes), fixed before any verifier exists, that collects and runs offline within \dv{screen.s10.minutes} minutes.
\item B1 is the subset's tests at the addressing commit that pass on the reference, or unavailable.
\item Repository code ran only in containers without network, and every rejection is ledgered with its reason (D59).
\item Building stops at \dv{design.family.cap} accepted tasks or once every family of both samples is processed
  (protocol~\S19.5); protocol v0.12 gives each sample a run of its own (\S20.2), and \texttt{\NATURALRUN} graded the
  first sample once building had stopped for it.
\end{itemize}

\paragraph{The task.} The initial request is the linked issue's title and body; the correction is the review comment
with the same reviewer's later replies before the addressing commit, plus its diff hunk; both are sanitised. The
pre-review state, the tree at the reviewed commit, is a held-out candidate (origin \texttt{natural\_pre\_review}); the
reference is the tree at the addressing commit. Tasks hold pseudonymised third-party text and live in a git-ignored
directory with a committed index of ids, hashes and licences (protocol~\S17.4); this paper prints no natural text,
repository name or identifier beyond the study's task ids. A review comment is not a chat correction, and the reviewer
need not be the person who asked for the change. We treat such a comment as the correction, a construct difference from
an in-session correction: this is D64's decision~1, a builder default that the owner has not reviewed (D65; first raised
in D47).

\paragraph{Replay of the first sample.} S10 of the first sample processed \nat{s10.first.processed} families (the
reserves NP-21 and NP-22 replaced NP-01 and NP-02 and were also rejected) and passed \nat{s10.first.passed}:
\nat{s11.ids} (D59; Table~\ref{tab:natural-s10}), whose B1 verifiers were frozen before anything else was generated for
them (D64). Their licences, as AIDev records them and as the LICENSE files at their bases state them, are in
Table~\ref{tab:natural-s10}. NP-03's snapshot is approximate, and those of NP-04, NP-05,
NP-18, NP-19 and NP-20 are exact (Appendix~\ref{app:construction}). The extension's \nat{fam.ext} families and
\nat{fam.ext.reserves} reserves (\nat{ni.ext.first} to \nat{ni.ext.last}) have no S10 outcome in the committed index:
its S10 step was stopped at a usage limit. The task directories that the stopped step left, NP-26 to NP-35, hold S10
outcomes that it wrote after the limit, which the index does not record, and they are to be built again from the start
(D66, D75), for \texttt{natural-pr-v2} (protocol~\S20.2).

\paragraph{Infrastructure attrition.} The first sample's S10 rejections are environmental or follow from S10's rules:
base tests that fail to collect (\nat{s10.first.reject.collect}), no resolvable dependency closure
(\nat{s10.first.reject.closure}), an environment over the size cap (\nat{s10.first.reject.toolarge}), an empty base-test
subset (\nat{s10.first.reject.emptysubset}), a fetch over the repository cap (\nat{s10.first.reject.fetchcap}) and the
merge rule (\nat{s10.first.reject.merge}). The rejections of NP-10, NP-16 and NP-17 depend on pip's resolver and on the
compiler's resolver limit of \dv{screen.s10.resolver.minutes} minutes, an implementation parameter rather than a
protocol rule; others follow from rules that cannot yield a runnable environment: path requirements outside the
workspace, generated version modules and submodules. Replay therefore favours small repositories with plain test
set-ups and fails on environments that cannot be pinned offline (D59).

\IfFileExists{generated/natural/natural_s10.tex}{\begin{table}[t]\centering\scriptsize\setlength{\tabcolsep}{4pt}
\caption{Stage S10 (replayability) of the first sample's 20 families and 5 reserves, from the committed task index. Fidelity: exact, or approximate with the number of approximations. Licence: the source repository's licence in AIDev; $^{\dagger}$: the LICENSE file at the base disagrees. Pins: the pinned dependency closure. Subset: test files in the base-test subset; base tests at the base: passed/failed/skipped/error. B1: the repository-test verifier is available. Reason: the S10 rejection category of the funnel table (no dependency closure: no resolvable dependency closure; fetch over the cap: fetch over the repository cap). Reserves are built only to replace a rejected family of their repository. The 23 families and reserves of the S9 extension (NP-26 to NP-48): \unfilled{pending: S10 of the S9 extension}.}\label{tab:natural-s10}
\begin{tabular}{llllrrrlllc}\toprule Task & Role & Fidelity & Licence & Pins & Subset & Base tests & B1 & S10 & Reason & S11 \\\midrule
NP-01 & family & approx.\ (1) & MIT & 24 & 1 & 0/0/0/1 & -- & rejected & base tests fail to collect & -- \\
NP-02 & family & approx.\ (1) & MIT & 95 & 10 & 0/0/0/0 & -- & rejected & base tests fail to collect & -- \\
NP-03 & family & approx.\ (1) & MIT & 73 & 14 & 954/10/2/0 & yes & passed & -- & yes \\
NP-04 & family & exact & MIT & 17 & 1 & 33/0/0/0 & yes & passed & -- & yes \\
NP-05 & family & exact & MIT & 17 & 16 & 354/5/0/0 & yes & passed & -- & yes \\
NP-06 & family & approx.\ (1) & Apache-2.0 & 92 & 1 & 0/0/0/1 & -- & rejected & base tests fail to collect & -- \\
NP-07 & family & -- & -- & -- & -- & -- & -- & rejected & empty base-test subset & -- \\
NP-08 & family & -- & -- & -- & -- & -- & -- & rejected & fetch over the cap & -- \\
NP-09 & family & approx.\ (1) & Apache-2.0 & 146 & 2 & 0/0/0/2 & -- & rejected & base tests fail to collect & -- \\
NP-10 & family & approx.\ (1) & Apache-2.0 & -- & 22 & -- & -- & rejected & no dependency closure & -- \\
NP-11 & family & approx.\ (2) & BSD-3-Clause & 80 & 42 & 0/0/0/23 & -- & rejected & base tests fail to collect & -- \\
NP-12 & family & -- & -- & -- & -- & -- & -- & rejected & merge rule & -- \\
NP-13 & family & exact & MIT & -- & 1 & -- & -- & rejected & no dependency closure & -- \\
NP-14 & family & exact & MIT & 52 & 1 & 0/0/0/13 & -- & rejected & base tests fail to collect & -- \\
NP-15 & family & approx.\ (1) & Apache-2.0 & 71 & 1 & -- & -- & rejected & environment too large & -- \\
NP-16 & family & approx.\ (1) & MIT & -- & 17 & -- & -- & rejected & no dependency closure & -- \\
NP-17 & family & approx.\ (1) & MIT & -- & 1 & -- & -- & rejected & no dependency closure & -- \\
NP-18 & family & exact & MIT & 144 & 1 & 24/0/0/0 & yes & passed & -- & yes \\
NP-19 & family & exact & Apache-2.0 & 106 & 2 & 1/5/0/3 & yes & passed & -- & yes \\
NP-20 & family & exact & MIT & 71 & 3 & 48/0/0/0 & yes & passed & -- & yes \\
NP-21 & reserve & approx.\ (1) & MIT & 123 & 50 & 0/0/0/0 & -- & rejected & base tests fail to collect & -- \\
NP-22 & reserve & approx.\ (1) & MIT & 32 & 2 & 0/0/0/0 & -- & rejected & base tests fail to collect & -- \\
NP-23 & reserve & -- & -- & -- & -- & -- & -- & not built & -- & -- \\
NP-24 & reserve & -- & -- & -- & -- & -- & -- & not built & -- & -- \\
NP-25 & reserve & -- & -- & -- & -- & -- & -- & not built & -- & -- \\
\bottomrule\end{tabular}\end{table}
}{\unfilled{S10 table}}
\IfFileExists{generated/natural/trace_commons.tex}{\begin{table}[t]\centering\footnotesize
\caption{The Trace Commons screen (population natural:chat): 30 donated sessions screened, 30 rejected, 0 eligible. Reasons are multi-valued; each session counts once per category. Indented rows: sessions per recorded detail; a session can carry several details of one category, so these rows can sum to more than the category. No task was built from this source.}\label{tab:trace-commons}
\begin{tabular}{lr}\toprule Rejection reason & Sessions \\\midrule
no recoverable base commit & 29 \\
\quad \texttt{none} & 26 \\
\quad \texttt{commit\_seen\_repository\_unidentified} & 2 \\
\quad \texttt{repository\_identified\_no\_commit} & 1 \\
replay: unreconstructable without a base tree & 29 \\
no code licence identified & 25 \\
no correction candidate under the correction-cue rule & 15 \\
unsuitable task & 8 \\
\quad no file changed by the main agent & 5 \\
\quad only documentation, configuration or data files changed (non-testable objective) & 3 \\
\quad the first request is an operation (install, deploy, commit) & 2 \\
unsupported transcript format & 2 \\
environment failure & 1 \\
\bottomrule\end{tabular}\end{table}
}{\unfilled{Trace Commons table}}

\paragraph{Construction.} For the accepted families the recipe was fixed and logged before any natural agent ran (D64,
D65): \dv{natural.templates.w} prompt templates with their hashes, the order of work with its freeze points, and the
rendered inputs and proposer prompts of each family. Each family's held-out pool holds the pre-review state,
\dv{natural.pool.alternatives.w} valid alternatives and \dv{natural.pool.nearmisses.w} near misses, at least
\dv{natural.pool.correction.w} of which target the correction (protocol~\S15.6). The phases, each ending with a committed record (D64, D65):
\begin{itemize}
\item P0: the task's validation, its PR record and image, and B1 frozen;
\item P1: the contract proposers (initial-only and correction-aware) and the host's build of their contracts;
\item P2: the generators of C1, C3, B2 and B3, and the development attacks and alternatives; C1, C3 and B2 frozen;
\item P3: the audits of C2 and C4 with their revisers, and B3's per-test selection, each frozen;
\item P4: the reference contract, drafted after every generated verifier is frozen;
\item P5: the reference checks with their review loop, and the independent labeler and its revision, each frozen with a
  status record;
\item P6: the near-miss targets, then the held-out pool, each candidate in a workspace of its own;
\item P7: the mappers, the duplicate check and the audits.
\end{itemize}
The host's steps of a phase (the repository's commands and scripts) are run by \emph{host coordinators}: subagents of the
build session in configuration A, a coordinator per step, which write no generated content and whose context
reaches no generator (D65). The run is in Section~\ref{sec:natural-cases}, and its family paragraphs are in
Appendix~\ref{app:natural-cases}. Table~\ref{tab:natural-vs-controlled} sets the arm beside the controlled study,
Table~\ref{tab:nr-construction} gives each family's construction facts, Table~\ref{tab:nr-labeler} its second
labeling, Table~\ref{tab:nr-sanity} each condition's outcome on its reference and on its no-op, and
Table~\ref{tab:nr-coverage} the coverage of its reference contract by the recovered contracts, with their
requirements that contradict it and those they mark superseded.

\begin{table}[t]\centering\scriptsize
\caption{How the natural arm differs from the controlled study (protocol~\S15.4 to \S15.6, \S20; decision log D64, D70).
$^{\ddagger}$Natural case study; caveats in Section~\ref{sec:natural-cases}.}\label{tab:natural-vs-controlled}
\begin{tabular}{@{}>{\raggedright\arraybackslash}p{0.17\linewidth}>{\raggedright\arraybackslash}p{0.38\linewidth}>{\raggedright\arraybackslash}p{0.39\linewidth}@{}}\toprule
 & Controlled study (\texttt{\FRESHRUN}, \texttt{\PILOTRUN}) & Natural arm (\texttt{\NATURALRUN}) \\\midrule
Source & a real agent session on a pinned commit, whose user turns the study's agent wrote and hashed beforehand: synthetic
feedback & an agent's pull request in AIDev v5 and its review thread \\
Correction & $U_2$, scripted, and relayed to the session agent with a short harness frame & a reviewer's comment on the
pull request, with the same reviewer's later replies; the reviewer need not have opened the issue \\
First request; pre-correction state & $U_1$; the session agent's workspace after $U_1$ & the linked issue; the tree the
reviewer commented on (the pre-review state) \\
Reference & the session agent's final state & the addressing commit: evidence about the correction, not an oracle \\
Reference contract & the study's intended requirements, hashed before the session and operationalised with quotes by
the study's agent & drafted at P4, after every generated verifier was frozen, by a configuration-A agent from the
rendered thread, the addressing diff and the reference diff \\
Held-out candidates & a constructed pool (pilot) and fresh attempts & the pre-review state (from S10), and
\dv{natural.pool.alternatives.w} alternatives and \dv{natural.pool.nearmisses.w} near misses generated after the family's
verifiers were frozen \\
Labels & D19 from the reference checks and the regression norm (and, for constructed candidates, the construction
intent); the independent labeler beside it & D19 with \dv{natural.d19.repeats.w} repeats; the source-evidence label of the
pre-review state (the label of protocol~\S15.6's first estimand); the independent labeler beside it
(Table~\ref{tab:label-rules}) \\
Audit & holds the reference valid & holds the reference valid, so C2 and C4 are audited against an unreliable oracle \\
B1; B3 & the reference state's repository tests, the session agent's added tests among them; no B3 & the base-test
subset's tests at the addressing commit, equal to D19's regression norm in NP-03, NP-05 and NP-20 and in part in NP-19;
B3, diff-derived, a secondary arm \\
Reporting & \fv{families.w} families, descriptive; the per-family table is primary & exploratory case studies: per-family
tables only, with no pooled rate, interval, p-value or total \\
\bottomrule\end{tabular}\end{table}

\IfFileExists{generated/natural_run/natural_run_construction.tex}{\begin{table}[t]\centering\scriptsize\setlength{\tabcolsep}{4pt}
\caption{Construction facts per family (decision log D64, D70), from the analysis's per-family flags, construction, reliability and cost blocks. Contracts: requirements, scored active ones and those introduced at the correction (the review comment or a reply; with none, the reference contract does not observe the correction, D64 decision 14). Checks: the scored checks of each registered verifier. Audits (audit:v1, D27): revisions and the result, with the codes of the last iteration's findings when it failed. B3: available, taken out (its per-test selection kept no test) or not registered. B1 by S10: available, with its checks. B1 equals D19's regression norm: yes, in part or no, with B1's files that are the norm's, of B1's, and the norm's count; where they are equal, wholly or in part, a D19-labelled false rejection of B1 is uninformative to that extent (D64, decision 3). Reference checks: the status record written when they were frozen, with its fix rounds, and their passes on the reference and the no-op in the run. Cost: the generation records of the family's generated artifacts, the run's grading time and the S10 build time. Terms: --: not in the analysis; n/a: recorded as null; not recorded: no record of it (codes not recorded: an audit record without its iterations; fix rounds not recorded: no integer for them); unrecognised status: a record with neither status; not known: no regression-norm overlay. Run \texttt{natural-pr-v1} in case-study mode, as exploratory case studies (protocol \S19.6, \S20.3): the run grades 6 accepted tasks of its own (its families accepted in the committed index), fewer than 10, and each run is reported on its own, so every number is per family. Nothing is pooled over the families: no pooled rate, no interval and no p-value. No total is given. $^{\ddagger}$Natural case study; caveats in Section~\ref{sec:natural-cases}.}\label{tab:nr-construction}
\begin{tabular}{p{0.2\linewidth}p{0.112\linewidth}p{0.112\linewidth}p{0.112\linewidth}p{0.112\linewidth}p{0.112\linewidth}p{0.112\linewidth}}\toprule Fact & NP-03 & NP-04 & NP-05 & NP-18 & NP-19 & NP-20 \\\midrule
\raggedright\hspace{0pt}Sample & \raggedright\hspace{0pt}first & \raggedright\hspace{0pt}first & \raggedright\hspace{0pt}first & \raggedright\hspace{0pt}first & \raggedright\hspace{0pt}first & \raggedright\hspace{0pt}first \tabularnewline
\raggedright\hspace{0pt}Reference contract: requirements, scored active, at the correction & \raggedright\hspace{0pt}9, 7, 2 & \raggedright\hspace{0pt}5, 4, 2 & \raggedright\hspace{0pt}13, 10, 2 & \raggedright\hspace{0pt}5, 4, 2 & \raggedright\hspace{0pt}5, 4, 2 & \raggedright\hspace{0pt}16, 12, 5 \tabularnewline
\raggedright\hspace{0pt}Recovered contracts, scored active: initial-only, full-session & \raggedright\hspace{0pt}4, 6 & \raggedright\hspace{0pt}3, 3 & \raggedright\hspace{0pt}10, 11 & \raggedright\hspace{0pt}5, 3 & \raggedright\hspace{0pt}4, 3 & \raggedright\hspace{0pt}3, 5 \tabularnewline\midrule
\raggedright\hspace{0pt}Checks: C1, C2, C3, C4 & \raggedright\hspace{0pt}11, 14, 14, 15 & \raggedright\hspace{0pt}11, 12, 14, 14 & \raggedright\hspace{0pt}14, 14, 14, 15 & \raggedright\hspace{0pt}12, 12, 11, 12 & \raggedright\hspace{0pt}10, 13, 11, 15 & \raggedright\hspace{0pt}12, 13, 13, 13 \tabularnewline
\raggedright\hspace{0pt}Checks: B2, B1, B3 & \raggedright\hspace{0pt}14, 954, 11 & \raggedright\hspace{0pt}11, 34, 14 & \raggedright\hspace{0pt}12, 354, 14 & \raggedright\hspace{0pt}12, 26, 13 & \raggedright\hspace{0pt}13, 1, 8 & \raggedright\hspace{0pt}12, 48, 11 \tabularnewline
\raggedright\hspace{0pt}Checks: reference checks, labeler & \raggedright\hspace{0pt}19, 17 & \raggedright\hspace{0pt}18, 9 & \raggedright\hspace{0pt}18, 17 & \raggedright\hspace{0pt}14, 10 & \raggedright\hspace{0pt}10, 7 & \raggedright\hspace{0pt}20, 20 \tabularnewline\midrule
\raggedright\hspace{0pt}C2 audit: revisions, result & \raggedright\hspace{0pt}2, passed & \raggedright\hspace{0pt}3, failed (false rejection dev) & \raggedright\hspace{0pt}3, failed (false acceptance dev) & \raggedright\hspace{0pt}1, passed & \raggedright\hspace{0pt}2, passed & \raggedright\hspace{0pt}1, passed \tabularnewline
\raggedright\hspace{0pt}C4 audit: revisions, result & \raggedright\hspace{0pt}3, failed (reference not passing) & \raggedright\hspace{0pt}0, passed & \raggedright\hspace{0pt}3, failed (false acceptance dev) & \raggedright\hspace{0pt}1, passed & \raggedright\hspace{0pt}3, failed (reference not passing) & \raggedright\hspace{0pt}0, passed \tabularnewline
\raggedright\hspace{0pt}B3 & \raggedright\hspace{0pt}available & \raggedright\hspace{0pt}available & \raggedright\hspace{0pt}available & \raggedright\hspace{0pt}available & \raggedright\hspace{0pt}available & \raggedright\hspace{0pt}available \tabularnewline
\raggedright\hspace{0pt}B1 by S10: available (checks) & \raggedright\hspace{0pt}yes (954) & \raggedright\hspace{0pt}yes (34) & \raggedright\hspace{0pt}yes (354) & \raggedright\hspace{0pt}yes (26) & \raggedright\hspace{0pt}yes (1) & \raggedright\hspace{0pt}yes (48) \tabularnewline
\raggedright\hspace{0pt}B1 equals D19's regression norm & \raggedright\hspace{0pt}yes (14 of 14; norm 14) & \raggedright\hspace{0pt}no (0 of 1; norm 1) & \raggedright\hspace{0pt}yes (16 of 16; norm 16) & \raggedright\hspace{0pt}no (0 of 1; norm 1) & \raggedright\hspace{0pt}in part (1 of 2; norm 2) & \raggedright\hspace{0pt}yes (3 of 3; norm 3) \tabularnewline\midrule
\raggedright\hspace{0pt}Reference checks: status (fix rounds) & \raggedright\hspace{0pt}usable (0) & \raggedright\hspace{0pt}usable (0) & \raggedright\hspace{0pt}usable (2) & \raggedright\hspace{0pt}usable (0) & \raggedright\hspace{0pt}usable (1) & \raggedright\hspace{0pt}usable (0) \tabularnewline
\raggedright\hspace{0pt}Reference checks in the run: reference, no-op passes & \raggedright\hspace{0pt}1/1, 0/1 & \raggedright\hspace{0pt}1/1, 0/1 & \raggedright\hspace{0pt}1/1, 0/1 & \raggedright\hspace{0pt}1/1, 0/1 & \raggedright\hspace{0pt}1/1, 0/1 & \raggedright\hspace{0pt}1/1, 0/1 \tabularnewline
\raggedright\hspace{0pt}Held-out near misses on the correction & \raggedright\hspace{0pt}2 of 3 & \raggedright\hspace{0pt}2 of 3 & \raggedright\hspace{0pt}2 of 3 & \raggedright\hspace{0pt}2 of 3 & \raggedright\hspace{0pt}2 of 3 & \raggedright\hspace{0pt}2 of 3 \tabularnewline\midrule
\raggedright\hspace{0pt}Generated artifacts with a generation record & \raggedright\hspace{0pt}26 of 26 & \raggedright\hspace{0pt}26 of 26 & \raggedright\hspace{0pt}26 of 26 & \raggedright\hspace{0pt}26 of 26 & \raggedright\hspace{0pt}26 of 26 & \raggedright\hspace{0pt}26 of 26 \tabularnewline
\raggedright\hspace{0pt}Agent turns, tool calls & \raggedright\hspace{0pt}1{,}085, 1{,}391 & \raggedright\hspace{0pt}741, 940 & \raggedright\hspace{0pt}967, 1{,}185 & \raggedright\hspace{0pt}712, 837 & \raggedright\hspace{0pt}1{,}247, 1{,}532 & \raggedright\hspace{0pt}934, 1{,}183 \tabularnewline
\raggedright\hspace{0pt}Agent wall time (s) & \raggedright\hspace{0pt}34{,}949 & \raggedright\hspace{0pt}24{,}988 & \raggedright\hspace{0pt}31{,}106 & \raggedright\hspace{0pt}23{,}735 & \raggedright\hspace{0pt}55{,}541 & \raggedright\hspace{0pt}24{,}111 \tabularnewline
\raggedright\hspace{0pt}Grading wall time (s) & \raggedright\hspace{0pt}1{,}111 & \raggedright\hspace{0pt}422 & \raggedright\hspace{0pt}871 & \raggedright\hspace{0pt}5{,}733 & \raggedright\hspace{0pt}6{,}147 & \raggedright\hspace{0pt}1{,}310 \tabularnewline
\raggedright\hspace{0pt}S10 build time (s) & \raggedright\hspace{0pt}24 & \raggedright\hspace{0pt}7 & \raggedright\hspace{0pt}38 & \raggedright\hspace{0pt}199 & \raggedright\hspace{0pt}113 & \raggedright\hspace{0pt}67 \tabularnewline
\bottomrule\end{tabular}\end{table}
}{\unfilled{construction table (make\_tables.py -{}-natural-run)}}
\IfFileExists{generated/natural_run/natural_run_labeler.tex}{\begin{table}[t]\centering\footnotesize
\caption{The second labeling per family: the independent labeler (protocol \S15.5), a model proposal like D19. \textbf{Both labelings are model proposals; no human labels exist.} Status: the labeler's status record, with its round (one as written, two after its one revision); a labeler that is not usable is left out of the second labeling and reported. Labels: the family's held-out candidates under the second labeling (not graded: left out, or not graded by the labeler). Agreement with D19, each distinct graded patch once: agree/both decide (units), where both decide a unit both label valid or invalid; apart for the pre-review states, which existed before the labeler was written, and for the held-out pool. The pool's order: how many of its candidates the generation records show to have been generated after the labeler was frozen, of all, with those not after it and those whose order is not shown (no frozen time, or no generation start). Terms: not recorded: no status record; graded: in the second labeling; left out: its labeler is not usable; --: not in the analysis; n/a: recorded as null. Run \texttt{natural-pr-v1} in case-study mode, as exploratory case studies (protocol \S19.6, \S20.3): the run grades 6 accepted tasks of its own (its families accepted in the committed index), fewer than 10, and each run is reported on its own, so every number is per family. Nothing is pooled over the families: no pooled rate, no interval and no p-value. No agreement is pooled over the families. $^{\ddagger}$Natural case study; caveats in Section~\ref{sec:natural-cases}.}\label{tab:nr-labeler}
\begin{tabular}{p{0.25\linewidth}p{0.093\linewidth}p{0.093\linewidth}p{0.093\linewidth}p{0.093\linewidth}p{0.093\linewidth}p{0.093\linewidth}}\toprule Second labeling & NP-03 & NP-04 & NP-05 & NP-18 & NP-19 & NP-20 \\\midrule
\raggedright\hspace{0pt}Labeler status (round) & \raggedright\hspace{0pt}unusable (2) & \raggedright\hspace{0pt}usable (1) & \raggedright\hspace{0pt}usable (2) & \raggedright\hspace{0pt}usable (1) & \raggedright\hspace{0pt}unusable (2) & \raggedright\hspace{0pt}usable (2) \tabularnewline
\raggedright\hspace{0pt}In the second labeling & \raggedright\hspace{0pt}left out & \raggedright\hspace{0pt}graded & \raggedright\hspace{0pt}graded & \raggedright\hspace{0pt}graded & \raggedright\hspace{0pt}left out & \raggedright\hspace{0pt}graded \tabularnewline\midrule
\raggedright\hspace{0pt}Held-out candidates labelled valid & \raggedright\hspace{0pt}0 & \raggedright\hspace{0pt}2 & \raggedright\hspace{0pt}3 & \raggedright\hspace{0pt}2 & \raggedright\hspace{0pt}0 & \raggedright\hspace{0pt}3 \tabularnewline
\raggedright\hspace{0pt}Held-out candidates labelled invalid & \raggedright\hspace{0pt}0 & \raggedright\hspace{0pt}4 & \raggedright\hspace{0pt}3 & \raggedright\hspace{0pt}4 & \raggedright\hspace{0pt}0 & \raggedright\hspace{0pt}3 \tabularnewline
\raggedright\hspace{0pt}Held-out candidates labelled ambiguous & \raggedright\hspace{0pt}0 & \raggedright\hspace{0pt}0 & \raggedright\hspace{0pt}0 & \raggedright\hspace{0pt}0 & \raggedright\hspace{0pt}0 & \raggedright\hspace{0pt}0 \tabularnewline
\raggedright\hspace{0pt}Held-out candidates not graded & \raggedright\hspace{0pt}6 & \raggedright\hspace{0pt}0 & \raggedright\hspace{0pt}0 & \raggedright\hspace{0pt}0 & \raggedright\hspace{0pt}6 & \raggedright\hspace{0pt}0 \tabularnewline\midrule
\raggedright\hspace{0pt}Agreement with D19, pre-review states & \raggedright\hspace{0pt}0/0 (1) & \raggedright\hspace{0pt}1/1 (1) & \raggedright\hspace{0pt}1/1 (1) & \raggedright\hspace{0pt}1/1 (1) & \raggedright\hspace{0pt}0/0 (1) & \raggedright\hspace{0pt}1/1 (1) \tabularnewline
\raggedright\hspace{0pt}Agreement with D19, held-out pool & \raggedright\hspace{0pt}0/0 (5) & \raggedright\hspace{0pt}4/4 (5) & \raggedright\hspace{0pt}4/4 (5) & \raggedright\hspace{0pt}5/5 (5) & \raggedright\hspace{0pt}0/0 (5) & \raggedright\hspace{0pt}4/4 (5) \tabularnewline
\raggedright\hspace{0pt}Pool generated after the labeler was frozen & \raggedright\hspace{0pt}5 of 5 & \raggedright\hspace{0pt}5 of 5 & \raggedright\hspace{0pt}5 of 5 & \raggedright\hspace{0pt}5 of 5 & \raggedright\hspace{0pt}5 of 5 & \raggedright\hspace{0pt}5 of 5 \tabularnewline
\bottomrule\end{tabular}\end{table}
}{\unfilled{second-labeling table (make\_tables.py -{}-natural-run)}}
\IfFileExists{generated/natural_run/natural_run_sanity.tex}{\begin{table}[t]\centering\footnotesize
\caption{The sanity grading per family: each condition's outcome on the family's reference solution and on its no-op, the unchanged base, in the run, as reference, no-op: \checkmark{} accepted, $\times$ rejected, with a rejection that is not a task failure (an infrastructure outcome, or an unresolved pair whose repeats disagree, protocol \S15.3) named in brackets. The rates leave both out (protocol \S15.7). The secondary oracle-gated view keeps a condition's cell only where the condition accepts the reference, and Tables~\ref{tab:nr-source} and~\ref{tab:nr-fa-fr} mark with $^{\dagger}$ the cells it drops. Terms: --: the run did not grade the condition for the family; not graded: the run did not grade that candidate with the condition. Run \texttt{natural-pr-v1} in case-study mode, as exploratory case studies (protocol \S19.6, \S20.3): the run grades 6 accepted tasks of its own (its families accepted in the committed index), fewer than 10, and each run is reported on its own, so every number is per family. Nothing is pooled over the families: no pooled rate, no interval and no p-value. $^{\ddagger}$Natural case study; caveats in Section~\ref{sec:natural-cases}.}\label{tab:nr-sanity}
\begin{tabular}{p{0.25\linewidth}p{0.093\linewidth}p{0.093\linewidth}p{0.093\linewidth}p{0.093\linewidth}p{0.093\linewidth}p{0.093\linewidth}}\toprule Condition & NP-03 & NP-04 & NP-05 & NP-18 & NP-19 & NP-20 \\\midrule
\raggedright\hspace{0pt}C1 initial-only, plain & \raggedright\hspace{0pt}$\times$, $\times$ & \raggedright\hspace{0pt}\checkmark, $\times$ & \raggedright\hspace{0pt}$\times$, $\times$ & \raggedright\hspace{0pt}$\times$, $\times$ & \raggedright\hspace{0pt}$\times$, $\times$ & \raggedright\hspace{0pt}$\times$, $\times$ \tabularnewline
\raggedright\hspace{0pt}C2 initial-only, audited & \raggedright\hspace{0pt}\checkmark, $\times$ & \raggedright\hspace{0pt}\checkmark, $\times$ & \raggedright\hspace{0pt}\checkmark, $\times$ & \raggedright\hspace{0pt}\checkmark, $\times$ & \raggedright\hspace{0pt}\checkmark, $\times$ & \raggedright\hspace{0pt}\checkmark, $\times$ \tabularnewline
\raggedright\hspace{0pt}C3 correction-aware, plain & \raggedright\hspace{0pt}$\times$, $\times$ & \raggedright\hspace{0pt}\checkmark, $\times$ & \raggedright\hspace{0pt}$\times$, $\times$ & \raggedright\hspace{0pt}$\times$, $\times$ & \raggedright\hspace{0pt}$\times$, $\times$ & \raggedright\hspace{0pt}\checkmark, $\times$ \tabularnewline
\raggedright\hspace{0pt}C4 correction-aware, audited & \raggedright\hspace{0pt}$\times$, $\times$ & \raggedright\hspace{0pt}\checkmark, $\times$ & \raggedright\hspace{0pt}\checkmark, $\times$ & \raggedright\hspace{0pt}\checkmark, $\times$ & \raggedright\hspace{0pt}$\times$, $\times$ & \raggedright\hspace{0pt}\checkmark, $\times$ \tabularnewline
\raggedright\hspace{0pt}B2 session-LLM baseline & \raggedright\hspace{0pt}\checkmark, $\times$ & \raggedright\hspace{0pt}\checkmark, $\times$ & \raggedright\hspace{0pt}\checkmark, $\times$ & \raggedright\hspace{0pt}\checkmark, $\times$ & \raggedright\hspace{0pt}\checkmark, $\times$ & \raggedright\hspace{0pt}\checkmark, $\times$ \tabularnewline
\raggedright\hspace{0pt}B1 repository tests & \raggedright\hspace{0pt}\checkmark, \checkmark & \raggedright\hspace{0pt}\checkmark, $\times$ & \raggedright\hspace{0pt}\checkmark, \checkmark & \raggedright\hspace{0pt}\checkmark, $\times$ & \raggedright\hspace{0pt}\checkmark, \checkmark & \raggedright\hspace{0pt}\checkmark, \checkmark \tabularnewline\midrule
\raggedright\hspace{0pt}B3 diff-derived (secondary) & \raggedright\hspace{0pt}\checkmark, $\times$ & \raggedright\hspace{0pt}\checkmark, $\times$ & \raggedright\hspace{0pt}\checkmark, $\times$ & \raggedright\hspace{0pt}\checkmark, $\times$ & \raggedright\hspace{0pt}\checkmark, $\times$ & \raggedright\hspace{0pt}\checkmark, $\times$ \tabularnewline
\bottomrule\end{tabular}\end{table}
}{\unfilled{sanity table (make\_tables.py -{}-natural-run)}}
\IfFileExists{generated/natural_run/natural_run_coverage.tex}{\begin{table}[t]\centering\footnotesize
\caption{Coverage of each family's reference contract by its two recovered contracts, the initial-only and the correction-aware one, under each of the two model mappers, whose mappings are proposals: the reference contract's scored active requirements that a requirement of the recovered contract maps to without contradicting it, of all. Correction requirements: the reference contract's scored active requirements introduced at the correction (the review comment or a reply; decision log D64, decision 14), with, below, those that each recovered contract does not cover. Contradicting: the recovered contract's scored active requirements that the mapper judges to contradict a requirement of the reference contract, of all its scored active requirements. Marked superseded: the recovered contract's requirements whose status is superseded and which it does not mark unscored; they are not scored. Terms: --: not in the analysis (in the rows of a mapper: that mapper did not map that contract); none: no such requirement. Run \texttt{natural-pr-v1} in case-study mode, as exploratory case studies (protocol \S19.6, \S20.3): the run grades 6 accepted tasks of its own (its families accepted in the committed index), fewer than 10, and each run is reported on its own, so every number is per family. Nothing is pooled over the families: no pooled rate, no interval and no p-value. $^{\ddagger}$Natural case study; caveats in Section~\ref{sec:natural-cases}.}\label{tab:nr-coverage}
\begin{tabular}{p{0.25\linewidth}p{0.093\linewidth}p{0.093\linewidth}p{0.093\linewidth}p{0.093\linewidth}p{0.093\linewidth}p{0.093\linewidth}}\toprule Coverage & NP-03 & NP-04 & NP-05 & NP-18 & NP-19 & NP-20 \\\midrule
\raggedright\hspace{0pt}Correction requirements of the reference contract & \raggedright\hspace{0pt}R8, R9 & \raggedright\hspace{0pt}R4, R5 & \raggedright\hspace{0pt}R7, R8 & \raggedright\hspace{0pt}R2, R3 & \raggedright\hspace{0pt}R3, R4 & \raggedright\hspace{0pt}R4, R5, R6, R11, R13 \tabularnewline\midrule
\raggedright\hspace{0pt}Coverage, first mapper: initial-only, correction-aware & \raggedright\hspace{0pt}2/7, 5/7 & \raggedright\hspace{0pt}2/4, 3/4 & \raggedright\hspace{0pt}8/10, 7/10 & \raggedright\hspace{0pt}4/4, 3/4 & \raggedright\hspace{0pt}2/4, 3/4 & \raggedright\hspace{0pt}4/12, 6/12 \tabularnewline
\raggedright\hspace{0pt}Coverage, second mapper: initial-only, correction-aware & \raggedright\hspace{0pt}2/7, 5/7 & \raggedright\hspace{0pt}2/4, 3/4 & \raggedright\hspace{0pt}8/10, 7/10 & \raggedright\hspace{0pt}4/4, 3/4 & \raggedright\hspace{0pt}2/4, 3/4 & \raggedright\hspace{0pt}3/12, 6/12 \tabularnewline\midrule
\raggedright\hspace{0pt}Correction requirements not covered, first mapper: initial-only; correction-aware & \raggedright\hspace{0pt}R8, R9; none & \raggedright\hspace{0pt}R4, R5; R5 & \raggedright\hspace{0pt}none; none & \raggedright\hspace{0pt}none; none & \raggedright\hspace{0pt}R3, R4; R4 & \raggedright\hspace{0pt}R5, R6, R11, R13; R11 \tabularnewline
\raggedright\hspace{0pt}Correction requirements not covered, second mapper: initial-only; correction-aware & \raggedright\hspace{0pt}R8, R9; none & \raggedright\hspace{0pt}R4, R5; R5 & \raggedright\hspace{0pt}none; none & \raggedright\hspace{0pt}none; none & \raggedright\hspace{0pt}R3, R4; R4 & \raggedright\hspace{0pt}R4, R5, R6, R11, R13; R11 \tabularnewline\midrule
\raggedright\hspace{0pt}Contradicting the reference contract, first mapper: initial-only, correction-aware & \raggedright\hspace{0pt}0/4, 0/6 & \raggedright\hspace{0pt}0/3, 0/3 & \raggedright\hspace{0pt}2/10, 3/11 & \raggedright\hspace{0pt}0/5, 0/3 & \raggedright\hspace{0pt}0/4, 0/3 & \raggedright\hspace{0pt}1/3, 1/5 \tabularnewline
\raggedright\hspace{0pt}Contradicting the reference contract, second mapper: initial-only, correction-aware & \raggedright\hspace{0pt}0/4, 0/6 & \raggedright\hspace{0pt}0/3, 0/3 & \raggedright\hspace{0pt}2/10, 3/11 & \raggedright\hspace{0pt}0/5, 0/3 & \raggedright\hspace{0pt}0/4, 0/3 & \raggedright\hspace{0pt}1/3, 1/5 \tabularnewline
\raggedright\hspace{0pt}Marked superseded by the recovered contract: initial-only, correction-aware & \raggedright\hspace{0pt}0, 0 & \raggedright\hspace{0pt}0, 0 & \raggedright\hspace{0pt}0, 0 & \raggedright\hspace{0pt}0, 0 & \raggedright\hspace{0pt}0, 0 & \raggedright\hspace{0pt}0, 1 \tabularnewline
\bottomrule\end{tabular}\end{table}
}{\unfilled{coverage table (make\_tables.py -{}-natural-run)}}

% The natural section's per-family detail, moved here unchanged after the analysis (D80).
\clearpage% Appendix G's part of the natural case studies (Section sec:natural-cases): detail of the template of
% paper/sections/natural.tex, moved here unchanged after the analysis (D80; D82 for the run's dates and invocations), so
% that the main-text section stays short. It is part of the template: every rule of natural.tex's
% header holds for it, and tests/test_natural_template.py reads it with natural.tex.
% * Moved unchanged: the caveats on the build's process (from the list of caveats), the sentence that names the parts of
%   a family paragraph (from the paragraph on reading the tables), Table tab:nr-caveats, the six family paragraphs and
%   the cost per family (from the closing paragraph); in a later round (D82), the sentence on the run's start and end
%   times, invocations and code revision (from the Status paragraph), re-wrapped (line breaks only).
% * Written after the analysis (D80): the subsection's heading and first paragraph, and the titles of the paragraphs of
%   the caveats, of the parts and of the cost; in the later round (D82), the title of the paragraph of the run's dates
%   and invocations, and the first paragraph's mention of it, with its citations.
\subsection{The case studies in detail}
\label{app:natural-cases}
This subsection holds the parts of Section~\ref{sec:natural-cases} that were moved here unchanged after the
analysis (D80, D82): the run's dates and invocations from its status paragraph, its caveats on the build's process,
its table of per-family caveats, its family paragraphs and the cost per family. The caveats of
Section~\ref{sec:natural-cases} hold for all of it.

\paragraph{The run's dates and invocations.}
Its manifest records its start on \nrv{run.date} at \nrv{run.started} UTC, its end at \nrv{run.finished} UTC, and
\nrv{run.invocations} invocations at code revision \texttt{\nrv{run.revision}}; D76 records why the first invocation was
stopped, within seconds and before any pair was graded, and that the run was then started again.

\paragraph{Caveats on the build's process.}
\begin{itemize}
\item D19 graded the frozen reference checks and the regression norm on every held-out candidate as the run's first
  step, once building had stopped for the first sample: a builder reading of protocol~\S19.7 and \S20.2, pending the
  owner's review (D73, D74, D76).
\item Some reference-check writers, labelers and their revisions, reviews, alternatives, near misses and fix rounds of P5
  and P6 had in their context the titles of the build session's task-list items (about the paper's tables and review
  fixes), and so did some host coordinators, which write no generated content (D65, D72). No title names a family, requirement, check, target, label or candidate; a search of
  each exposed generator run's text and tool inputs for the titles' distinctive terms found none; and nothing was
  regenerated, a builder decision not yet reviewed by a person (D72).
\item Some runs worked outside the recipe: they started the task image themselves, wrote temporary files outside their
  workspaces, or queried the container engine read-only. Each was assessed and none is excluded (D75); the family
  paragraphs name those that D75 names.
\item The templates whose sources carry an owner-context line leave it out, a builder decision pending the owner's
  review (D64, deviation~(i) and decision~17). The drafter of the reference contract reads the reference diff beyond the
  inputs that protocol~\S15.6 names (D64, deviation~(k)).
\item The harness context. The runs of P1 and of P5 carried, in their first message, a harness entry relaying the
  owner's latest message (about the usage budget at P1, about finishing the study at P5), which names no task, input or
  output and which the harness ranks above the computed task where they conflict. It could still nudge an agent
  toward less or more work. The decision log records which runs carried it, and nothing was regenerated for it, a
  builder decision pending the owner's review (D67, D68, D77). The agent tool offers no per-call tool restriction, so each natural agent had its
  agent type's tools, as the controlled agents of the same source workflows had, and the natural-agent audit read every
  network use and tool call (D64, ``Agents'' and decision~11; D65). An environment block naming the study repository as
  the working directory was in every natural agent's context, as in the fresh attempts (D65; protocol~\S18.2).
\end{itemize}

\paragraph{The parts of a family paragraph.}
Each family paragraph gives, in this order, (a)~its construction facts, (b)~its pre-review state under the
source-evidence label, (c)~its held-out candidates under D19, (d)~its second labeling, (e)~its unresolved and
infrastructure pairs, and (f)~its own caveats.

\begin{table}[t]\centering\scriptsize\setlength{\tabcolsep}{4pt}
\caption{Per-family caveats of \texttt{\NATURALRUN}, each cell a value of the run's analysis (\texttt{make\_tables.py
-{}-natural-run}; decision log D64, D70, D72, D75, D76). Reference checks: the status record written when they were
frozen, with its fix rounds. Labeler: its status record, and whether the second labeling grades the family or leaves it
out. Oracle-gated view: kept where the verifier passes the family's reference solution in the run, dropped otherwise
(protocol~\S15.7). B1: whether it equals D19's regression norm (yes, in part or no), with its checks. Near misses: those
that target a correction requirement, of all. D19 ambiguous: the candidates D19 labels ambiguous, left out of false
acceptance and false rejection. Terms: not recorded: no status record; unrecognised status: a record with neither
status; not known: no regression-norm overlay; graded: in the second labeling; left out: the family's labeler is not
usable; not graded: the run did not grade the family's labeler; none: no such candidate. $^{\ddagger}$Natural case
study; caveats in Section~\ref{sec:natural-cases}.}\label{tab:nr-caveats}
\begin{tabular}{p{0.2\linewidth}p{0.112\linewidth}p{0.112\linewidth}p{0.112\linewidth}p{0.112\linewidth}p{0.112\linewidth}p{0.112\linewidth}}
\toprule
\nrcaveat{Caveat} & \nrcaveat{NP-03} & \nrcaveat{NP-04} & \nrcaveat{NP-05} & \nrcaveat{NP-18} & \nrcaveat{NP-19} &
\nrcaveat{NP-20} \tabularnewline\midrule
\nrcaveat{Reference checks: status (fix rounds)}
& \nrcaveat{\nrv{refchecks.np03.status} (\nrv{refchecks.np03.rounds})}
& \nrcaveat{\nrv{refchecks.np04.status} (\nrv{refchecks.np04.rounds})}
& \nrcaveat{\nrv{refchecks.np05.status} (\nrv{refchecks.np05.rounds})}
& \nrcaveat{\nrv{refchecks.np18.status} (\nrv{refchecks.np18.rounds})}
& \nrcaveat{\nrv{refchecks.np19.status} (\nrv{refchecks.np19.rounds})}
& \nrcaveat{\nrv{refchecks.np20.status} (\nrv{refchecks.np20.rounds})} \tabularnewline
\nrcaveat{Labeler: status; second labeling}
& \nrcaveat{\nrv{il.np03.status}; \nrv{il.np03.second}}
& \nrcaveat{\nrv{il.np04.status}; \nrv{il.np04.second}}
& \nrcaveat{\nrv{il.np05.status}; \nrv{il.np05.second}}
& \nrcaveat{\nrv{il.np18.status}; \nrv{il.np18.second}}
& \nrcaveat{\nrv{il.np19.status}; \nrv{il.np19.second}}
& \nrcaveat{\nrv{il.np20.status}; \nrv{il.np20.second}} \tabularnewline
\nrcaveat{Oracle-gated view: C1, C3}
& \nrcaveat{\nrv{og.c1.np03}, \nrv{og.c3.np03}}
& \nrcaveat{\nrv{og.c1.np04}, \nrv{og.c3.np04}}
& \nrcaveat{\nrv{og.c1.np05}, \nrv{og.c3.np05}}
& \nrcaveat{\nrv{og.c1.np18}, \nrv{og.c3.np18}}
& \nrcaveat{\nrv{og.c1.np19}, \nrv{og.c3.np19}}
& \nrcaveat{\nrv{og.c1.np20}, \nrv{og.c3.np20}} \tabularnewline
\nrcaveat{B1: equals the norm (checks)}
& \nrcaveat{\nrv{b1.np03.norm} (\nrv{b1.np03.checks})}
& \nrcaveat{\nrv{b1.np04.norm} (\nrv{b1.np04.checks})}
& \nrcaveat{\nrv{b1.np05.norm} (\nrv{b1.np05.checks})}
& \nrcaveat{\nrv{b1.np18.norm} (\nrv{b1.np18.checks})}
& \nrcaveat{\nrv{b1.np19.norm} (\nrv{b1.np19.checks})}
& \nrcaveat{\nrv{b1.np20.norm} (\nrv{b1.np20.checks})} \tabularnewline
\nrcaveat{Near misses on the correction}
& \nrcaveat{\nrv{nm.np03.correction} of \nrv{nm.np03.all}}
& \nrcaveat{\nrv{nm.np04.correction} of \nrv{nm.np04.all}}
& \nrcaveat{\nrv{nm.np05.correction} of \nrv{nm.np05.all}}
& \nrcaveat{\nrv{nm.np18.correction} of \nrv{nm.np18.all}}
& \nrcaveat{\nrv{nm.np19.correction} of \nrv{nm.np19.all}}
& \nrcaveat{\nrv{nm.np20.correction} of \nrv{nm.np20.all}} \tabularnewline
\nrcaveat{D19 ambiguous}
& \nrcaveat{\nrv{d19.np03.ambiguous.ids}}
& \nrcaveat{\nrv{d19.np04.ambiguous.ids}}
& \nrcaveat{\nrv{d19.np05.ambiguous.ids}}
& \nrcaveat{\nrv{d19.np18.ambiguous.ids}}
& \nrcaveat{\nrv{d19.np19.ambiguous.ids}}
& \nrcaveat{\nrv{d19.np20.ambiguous.ids}} \tabularnewline
\bottomrule
\end{tabular}
\end{table}

\paragraph{NP-03.}
NP-03 has no second labeling: its independent labeler is unusable, since some of its checks state observables that the
reference does not meet (D72).
(a)~Construction (D64, D70, D72). Reference contract (requirements; scored and active; at the correction):
\nrv{ref.np03.requirements}; \nrv{ref.np03.scored}; \nrv{ref.np03.correction}. Reference checks:
\nrv{refchecks.np03.status}, fix rounds \nrv{refchecks.np03.rounds}. Labeler: \nrv{il.np03.status}; second labeling:
\nrv{il.np03.second}. B1 by S10 (available; checks; equal to D19's regression norm): \nrv{b1.np03.available};
\nrv{b1.np03.checks}; \nrv{b1.np03.norm}. Near misses on the correction: \nrv{nm.np03.correction} of
\nrv{nm.np03.all}. D19's labels of the held-out candidates (valid; invalid; ambiguous; no proposal):
\nrv{d19.np03.valid}; \nrv{d19.np03.invalid}; \nrv{d19.np03.ambiguous}; \nrv{d19.np03.unlabeled}; the ambiguous ones
with D19's reason: \nrv{d19.np03.ambiguous.reasons}.
(b)~The pre-review state, invalid under the source-evidence label (protocol~\S15.6): C1 \nrv{se.c1.np03}, C3
\nrv{se.c3.np03}; C2 \nrv{se.c2.np03} and C4 \nrv{se.c4.np03}, audited against an unreliable oracle (D70); B2
\nrv{se.b2.np03}; B1 \nrv{se.b1.np03}; B3, the secondary arm, \nrv{se.b3.np03}. Its D19 label: \nrv{pre.np03.d19}; the
reference checks' outcome on it in the run: \nrv{pre.np03.refchecks}.
(c)~Under D19: C1 \nrv{c1.np03.fa} and \nrv{c1.np03.fr}, C3 \nrv{c3.np03.fa} and \nrv{c3.np03.fr}; in the oracle-gated
view C1's cell is \nrv{og.c1.np03} and C3's \nrv{og.c3.np03}. B1 equals D19's regression norm here, so its false
rejection, \nrv{b1.np03.fr}, is uninformative, and its verdicts on the alternatives are not blind: they were developed
against B1's own test files (D64, decision~3). No near-miss writer ran the repository's tests, the writer of eval-nm-3
included, which started the task image itself, as the family's audit records of P7 say (D75, D77).
(d)~Second labeling: none, since the labeler is unusable (D72).
(e)~Pairs of the run (graded; unresolved; infrastructure outcomes): \nrv{pairs.np03}; \nrv{pairs.np03.unresolved};
\nrv{pairs.np03.infrastructure}.
(f)~The reference-check writer found that the reference cannot meet the reference contract's observables R1, R8 and
R9, and wrote checks that hold either way there; the near misses target R8, R9 and R1, so D19's labels for them do not
rest on checks of the targeted observables as the contract states them (D72). The review of the reference checks was
ok, with coverage remarks on R2 and R5, so no fix round ran (D33, D72). The base snapshot is approximate (D59). At the
audit's first iteration the reference fails C1 and C3, and the audit of C4 failed after its last revision with the
reference not passing (D70). The writer of eval-nm-3 started the task image with no network, its own read-only
repository and the wrappers' isolation options (D75).

\paragraph{NP-04.}
(a)~Construction (D64, D70, D72). Reference contract (requirements; scored and active; at the correction):
\nrv{ref.np04.requirements}; \nrv{ref.np04.scored}; \nrv{ref.np04.correction}. Reference checks:
\nrv{refchecks.np04.status}, fix rounds \nrv{refchecks.np04.rounds}. Labeler: \nrv{il.np04.status}; second labeling:
\nrv{il.np04.second}. B1 by S10 (available; checks; equal to D19's regression norm): \nrv{b1.np04.available};
\nrv{b1.np04.checks}; \nrv{b1.np04.norm}. Near misses on the correction: \nrv{nm.np04.correction} of
\nrv{nm.np04.all}. D19's labels of the held-out candidates (valid; invalid; ambiguous; no proposal):
\nrv{d19.np04.valid}; \nrv{d19.np04.invalid}; \nrv{d19.np04.ambiguous}; \nrv{d19.np04.unlabeled}; the ambiguous ones
with D19's reason: \nrv{d19.np04.ambiguous.reasons}.
(b)~The pre-review state, invalid under the source-evidence label (protocol~\S15.6): C1 \nrv{se.c1.np04}, C3
\nrv{se.c3.np04}; C2 \nrv{se.c2.np04} and C4 \nrv{se.c4.np04}, audited against an unreliable oracle (D70); B2
\nrv{se.b2.np04}; B1 \nrv{se.b1.np04}; B3, the secondary arm, \nrv{se.b3.np04}. Its D19 label: \nrv{pre.np04.d19}; the
reference checks' outcome on it in the run: \nrv{pre.np04.refchecks}.
(c)~Under D19: C1 \nrv{c1.np04.fa} and \nrv{c1.np04.fr}, C3 \nrv{c3.np04.fa} and \nrv{c3.np04.fr}; in the oracle-gated
view C1's cell is \nrv{og.c1.np04} and C3's \nrv{og.c3.np04}. No near-miss writer ran the repository's tests, as the
family's audit records of P7 say (D77).
(d)~Second labeling, on the held-out pool (units both label alike, of the units both decide): \nrv{il.np04.pool.agree}
of \nrv{il.np04.pool.decided}.
(e)~Pairs of the run (graded; unresolved; infrastructure outcomes): \nrv{pairs.np04}; \nrv{pairs.np04.unresolved};
\nrv{pairs.np04.infrastructure}.
(f)~NP-04 shares its repository with NP-05 (D63). At the audit's first iteration the reference passes C1 and C3 (D70).
The audit of C2 failed after its last revision by rejecting the development alternative, which, by the revisers'
account, cannot be told apart from the reference on the requirement it targets (D70). The writer of eval-alt-1 wrote and
ran a helper outside its workspace; a review of the reference checks started the task image itself, its verdict rests
partly on those runs, and the audit of P7 and the strict re-audit missed a call it made (D75).

\paragraph{NP-05.}
(a)~Construction (D64, D70, D75). Reference contract (requirements; scored and active; at the correction):
\nrv{ref.np05.requirements}; \nrv{ref.np05.scored}; \nrv{ref.np05.correction}. Reference checks:
\nrv{refchecks.np05.status}, fix rounds \nrv{refchecks.np05.rounds}. Labeler: \nrv{il.np05.status}; second labeling:
\nrv{il.np05.second}. B1 by S10 (available; checks; equal to D19's regression norm): \nrv{b1.np05.available};
\nrv{b1.np05.checks}; \nrv{b1.np05.norm}. Near misses on the correction: \nrv{nm.np05.correction} of
\nrv{nm.np05.all}. D19's labels of the held-out candidates (valid; invalid; ambiguous; no proposal):
\nrv{d19.np05.valid}; \nrv{d19.np05.invalid}; \nrv{d19.np05.ambiguous}; \nrv{d19.np05.unlabeled}; the ambiguous ones
with D19's reason: \nrv{d19.np05.ambiguous.reasons}.
(b)~The pre-review state, invalid under the source-evidence label (protocol~\S15.6): C1 \nrv{se.c1.np05}, C3
\nrv{se.c3.np05}; C2 \nrv{se.c2.np05} and C4 \nrv{se.c4.np05}, audited against an unreliable oracle (D70); B2
\nrv{se.b2.np05}; B1 \nrv{se.b1.np05}; B3, the secondary arm, \nrv{se.b3.np05}. Its D19 label: \nrv{pre.np05.d19}; the
reference checks' outcome on it in the run: \nrv{pre.np05.refchecks}.
(c)~Under D19: C1 \nrv{c1.np05.fa} and \nrv{c1.np05.fr}, C3 \nrv{c3.np05.fa} and \nrv{c3.np05.fr}; in the oracle-gated
view C1's cell is \nrv{og.c1.np05} and C3's \nrv{og.c3.np05}. B1 equals D19's regression norm here, so its false
rejection, \nrv{b1.np05.fr}, is uninformative, and its verdicts on the alternatives are not blind: they were developed
against B1's own test files (D64, decision~3). No near-miss writer ran the repository's tests, as the family's audit
records of P7 say (D77).
(d)~Second labeling, on the held-out pool (units both label alike, of the units both decide): \nrv{il.np05.pool.agree}
of \nrv{il.np05.pool.decided}.
(e)~Pairs of the run (graded; unresolved; infrastructure outcomes): \nrv{pairs.np05}; \nrv{pairs.np05.unresolved};
\nrv{pairs.np05.infrastructure}.
(f)~NP-05 shares its repository with NP-04 (D63). A qualifying comment of the extension was made before the family's
correction, and the family keeps the correction it was built with, a builder default that departs from the
earliest-comment rule of protocol~\S15.6 (D63, D66). At the audit's first iteration the
reference fails C1 and C3, and the audits of C2 and C4 failed after their last revision by accepting a development
attack that, by the revisers' account, cannot be told apart from the reference on the requirement it targets (D70). Its
reference checks were frozen after \nrv{refchecks.np05.rounds} fix rounds while the last review still asked for fixes,
which the review loop records without a further round (D33, D75). The writer of eval-nm-3 remarks that a strict reader
could also count it against R8; D19 decides its label (D75). The labeler's writer made a stub module outside its
workspace and then removed that directory, which also removed a stale stub of the controlled arm's tooling, whose
content is in those agents' transcripts (D75).

\paragraph{NP-18.}
(a)~Construction (D64, D70, D72). Reference contract (requirements; scored and active; at the correction):
\nrv{ref.np18.requirements}; \nrv{ref.np18.scored}; \nrv{ref.np18.correction}. Reference checks:
\nrv{refchecks.np18.status}, fix rounds \nrv{refchecks.np18.rounds}. Labeler: \nrv{il.np18.status}; second labeling:
\nrv{il.np18.second}. B1 by S10 (available; checks; equal to D19's regression norm): \nrv{b1.np18.available};
\nrv{b1.np18.checks}; \nrv{b1.np18.norm}. Near misses on the correction: \nrv{nm.np18.correction} of
\nrv{nm.np18.all}. D19's labels of the held-out candidates (valid; invalid; ambiguous; no proposal):
\nrv{d19.np18.valid}; \nrv{d19.np18.invalid}; \nrv{d19.np18.ambiguous}; \nrv{d19.np18.unlabeled}; the ambiguous ones
with D19's reason: \nrv{d19.np18.ambiguous.reasons}.
(b)~The pre-review state, invalid under the source-evidence label (protocol~\S15.6): C1 \nrv{se.c1.np18}, C3
\nrv{se.c3.np18}; C2 \nrv{se.c2.np18} and C4 \nrv{se.c4.np18}, audited against an unreliable oracle (D70); B2
\nrv{se.b2.np18}; B1 \nrv{se.b1.np18}; B3, the secondary arm, \nrv{se.b3.np18}. Its D19 label: \nrv{pre.np18.d19}; the
reference checks' outcome on it in the run: \nrv{pre.np18.refchecks}.
(c)~Under D19: C1 \nrv{c1.np18.fa} and \nrv{c1.np18.fr}, C3 \nrv{c3.np18.fa} and \nrv{c3.np18.fr}; in the oracle-gated
view C1's cell is \nrv{og.c1.np18} and C3's \nrv{og.c3.np18}. No near-miss writer ran the repository's tests, as the
family's audit records of P7 say (D77).
(d)~Second labeling, on the held-out pool (units both label alike, of the units both decide): \nrv{il.np18.pool.agree}
of \nrv{il.np18.pool.decided}.
(e)~Pairs of the run (graded; unresolved; infrastructure outcomes): \nrv{pairs.np18}; \nrv{pairs.np18.unresolved};
\nrv{pairs.np18.infrastructure}.
(f)~At the audit's first iteration the reference fails C1 and C3, and revisers relaxed checks that the reference fails,
so C2 and C4 can accept what the contract rules out (D70). The review of the reference checks was ok, with a coverage
remark on R4, so no fix round ran (D33, D72).

\paragraph{NP-19.}
NP-19 has no second labeling: its independent labeler is unusable after its revision, since some of its checks fail on
the reference and others pass on the no-op (D75).
(a)~Construction (D64, D70, D75). Reference contract (requirements; scored and active; at the correction):
\nrv{ref.np19.requirements}; \nrv{ref.np19.scored}; \nrv{ref.np19.correction}. Reference checks:
\nrv{refchecks.np19.status}, fix rounds \nrv{refchecks.np19.rounds}. Labeler: \nrv{il.np19.status}; second labeling:
\nrv{il.np19.second}. B1 by S10 (available; checks; equal to D19's regression norm): \nrv{b1.np19.available};
\nrv{b1.np19.checks}; \nrv{b1.np19.norm}. Near misses on the correction: \nrv{nm.np19.correction} of
\nrv{nm.np19.all}. D19's labels of the held-out candidates (valid; invalid; ambiguous; no proposal):
\nrv{d19.np19.valid}; \nrv{d19.np19.invalid}; \nrv{d19.np19.ambiguous}; \nrv{d19.np19.unlabeled}; the ambiguous ones
with D19's reason: \nrv{d19.np19.ambiguous.reasons}.
(b)~The pre-review state, invalid under the source-evidence label (protocol~\S15.6): C1 \nrv{se.c1.np19}, C3
\nrv{se.c3.np19}; C2 \nrv{se.c2.np19} and C4 \nrv{se.c4.np19}, audited against an unreliable oracle (D70); B2
\nrv{se.b2.np19}; B1 \nrv{se.b1.np19}; B3, the secondary arm, \nrv{se.b3.np19}. Its D19 label: \nrv{pre.np19.d19}; the
reference checks' outcome on it in the run: \nrv{pre.np19.refchecks}.
(c)~Under D19: C1 \nrv{c1.np19.fa} and \nrv{c1.np19.fr}, C3 \nrv{c3.np19.fa} and \nrv{c3.np19.fr}; in the oracle-gated
view C1's cell is \nrv{og.c1.np19} and C3's \nrv{og.c3.np19}. B1 equals D19's regression norm in part here, so its false
rejection, \nrv{b1.np19.fr}, is uninformative to that extent (D64, decision~3). No near-miss writer ran the repository's
tests, as the family's audit records of P7 say (D77).
(d)~Second labeling: none, since the labeler is unusable (D75).
(e)~Pairs of the run (graded; unresolved; infrastructure outcomes): \nrv{pairs.np19}; \nrv{pairs.np19.unresolved};
\nrv{pairs.np19.infrastructure}.
(f)~At the audit's first iteration the reference fails C1 and C3; the audit of C4 failed after its last revision with
the reference not passing, and that revision changed nothing; and the generators report that the fix does not take
effect in the task's environment (D70). The second review of the reference checks, after their fix round, was ok, with a
coverage remark on R1, so no further fix round ran (D33, D75). The near-miss writers report that the reference fails R1
in their mocked setups, while the frozen reference checks pass the reference in every repeat of their validation (D75).
Neither recovered contract covers the reference contract's R4 (D75). A fix round of the reference checks started the
task image itself, and the writer of B3 wrote a file outside its workspace (D75).

\paragraph{NP-20.}
(a)~Construction (D64, D70, D72). Reference contract (requirements; scored and active; at the correction):
\nrv{ref.np20.requirements}; \nrv{ref.np20.scored}; \nrv{ref.np20.correction}. Reference checks:
\nrv{refchecks.np20.status}, fix rounds \nrv{refchecks.np20.rounds}. Labeler: \nrv{il.np20.status}; second labeling:
\nrv{il.np20.second}. B1 by S10 (available; checks; equal to D19's regression norm): \nrv{b1.np20.available};
\nrv{b1.np20.checks}; \nrv{b1.np20.norm}. Near misses on the correction: \nrv{nm.np20.correction} of
\nrv{nm.np20.all}. D19's labels of the held-out candidates (valid; invalid; ambiguous; no proposal):
\nrv{d19.np20.valid}; \nrv{d19.np20.invalid}; \nrv{d19.np20.ambiguous}; \nrv{d19.np20.unlabeled}; the ambiguous ones
with D19's reason: \nrv{d19.np20.ambiguous.reasons}.
(b)~The pre-review state, invalid under the source-evidence label (protocol~\S15.6): C1 \nrv{se.c1.np20}, C3
\nrv{se.c3.np20}; C2 \nrv{se.c2.np20} and C4 \nrv{se.c4.np20}, audited against an unreliable oracle (D70); B2
\nrv{se.b2.np20}; B1 \nrv{se.b1.np20}; B3, the secondary arm, \nrv{se.b3.np20}. Its D19 label: \nrv{pre.np20.d19}; the
reference checks' outcome on it in the run: \nrv{pre.np20.refchecks}.
(c)~Under D19: C1 \nrv{c1.np20.fa} and \nrv{c1.np20.fr}, C3 \nrv{c3.np20.fa} and \nrv{c3.np20.fr}; in the oracle-gated
view C1's cell is \nrv{og.c1.np20} and C3's \nrv{og.c3.np20}. B1 equals D19's regression norm here, so its false
rejection, \nrv{b1.np20.fr}, is uninformative, and its verdicts on the alternatives are not blind: they were developed
against B1's own test files (D64, decision~3). The family's audit records of P7 do not say whether a near-miss writer ran
the repository's tests; a search of its near-miss writers' tool inputs for a fixed list of test-run commands
found none (D77).
(d)~Second labeling, on the held-out pool (units both label alike, of the units both decide): \nrv{il.np20.pool.agree}
of \nrv{il.np20.pool.decided}.
(e)~Pairs of the run (graded; unresolved; infrastructure outcomes): \nrv{pairs.np20}; \nrv{pairs.np20.unresolved};
\nrv{pairs.np20.infrastructure}.
(f)~The labeler is usable after its revision; as written it missed R15 (D72). The review of the reference checks was ok,
with a coverage remark on R11, so no fix round ran (D33, D72). At the audit's first iteration the reference fails C1 and
passes C3, and a reviser relaxed a check that the reference fails (D70). The writer of B2 wrote files of its error output
outside its workspace, in the host's temporary directory (D75).

\paragraph{Cost per family.}
Cost per family, from the generation records of the family's generated artifacts, the run's grading and its S10 record
(agent turns; tool calls; agent wall time, grading time and S10 build time in seconds; artifacts with a generation
record, of all; D64, ``Generation records''): NP-03 \nrv{cost.np03.turns}; \nrv{cost.np03.toolcalls};
\nrv{cost.np03.wall}, \nrv{cost.np03.grading} and \nrv{cost.np03.s10}; \nrv{cost.np03.recorded} of
\nrv{cost.np03.artifacts}. NP-04 \nrv{cost.np04.turns}; \nrv{cost.np04.toolcalls}; \nrv{cost.np04.wall},
\nrv{cost.np04.grading} and \nrv{cost.np04.s10}; \nrv{cost.np04.recorded} of \nrv{cost.np04.artifacts}. NP-05
\nrv{cost.np05.turns}; \nrv{cost.np05.toolcalls}; \nrv{cost.np05.wall}, \nrv{cost.np05.grading} and \nrv{cost.np05.s10};
\nrv{cost.np05.recorded} of \nrv{cost.np05.artifacts}. NP-18 \nrv{cost.np18.turns}; \nrv{cost.np18.toolcalls};
\nrv{cost.np18.wall}, \nrv{cost.np18.grading} and \nrv{cost.np18.s10}; \nrv{cost.np18.recorded} of
\nrv{cost.np18.artifacts}. NP-19 \nrv{cost.np19.turns}; \nrv{cost.np19.toolcalls}; \nrv{cost.np19.wall},
\nrv{cost.np19.grading} and \nrv{cost.np19.s10}; \nrv{cost.np19.recorded} of \nrv{cost.np19.artifacts}. NP-20
\nrv{cost.np20.turns}; \nrv{cost.np20.toolcalls}; \nrv{cost.np20.wall}, \nrv{cost.np20.grading} and \nrv{cost.np20.s10};
\nrv{cost.np20.recorded} of \nrv{cost.np20.artifacts}. This cost leaves out what no such record holds, as the run's
analysis lists it: the host's own steps, which the host coordinators run (D65), apart from the run's grading; agent runs
that leave no generation record, the mappers among them; stages S1 to S9; re-runs of S10 and the part of its compile
step that its record does not time; and image builds and wheel downloads (D64, ``Generation records''). So it is not the
whole cost of building a family: D68 reports the host coordinators' share of a phase. No total over the families is
given. The protocol's cost and time per accepted task (protocol~\S15.6), which the analysis computes over the families
that S10 processed, rejected ones included (D69), is a quantity over families and is not reported (protocol~\S19.6).

